\scanpage{51}
\deleted{b) L'ens des coorientations de}

b) \deleted{\unsure{$X$ et $\widetilde X$ sont les deux seuls bitriangles disjoints de $X$ et|$X$ et $\widetilde X$ sont les deux seuls bitriangles disjoints de $\bar X$}, \unsure{permutés} \unsure{d'où les trois} symétries}
\inserted{Considérons le groupe $\widetilde G$ d'autom de $X$, \deleted{\unsure{engendré}} engendré par les} symétries $\sigma_t$ $(t\in\mathrm{Tri}(\widetilde X))$. Alors

$\alpha$) $\widetilde G\longrightarrow\mathrm{Aut}_{\unsure{\mathrm{gr}|\mathrm{gr.}}}(\widetilde X)$ \deleted{est} \inserted{induit} un isom. \inserted{de $\widetilde G$} avec le groupe des autom \uline{intérieurs} de $\widetilde X$

$\beta$) Les orbites de $\widetilde G^{\unsure{\circ|0}}$ opérant sur $S-\widetilde X$ $(=X\amalg\bar X^{\unsure{+}})$ sont de cardinal $6$ -- donc il y \deleted{\unsure{en}} a trois telles orbites
\inserted{qui sont les réunions de la forme $\tau\cup\bar\tau$ (triade de l'hexagone \unsure{$H$})}
[On \unsure{a} en effet $\widetilde G\simeq\mathfrak S_{\unsure{P|E|\mathbb F}}\times\mathfrak S_Q$, où $\unsure{\bar P|\widetilde P|\mathbb F},\widetilde Q$ sont les deux ens de card. $3$ des triangles de $\widetilde X$ de l'une et l'autre espèce, \deleted{les orbites sont} $\mathfrak S_P$ et $\mathfrak S_Q$ sont \unsure{simplmt|simplemt|simplement} transitifs sur chacune des trois orbites]
\inserted{Donc l'ens des trois \unsure{orbites|calottes|culottes} s'identifie à l'ens $I$ des $3$ diagonales de $H$.}

c) \deleted{\unsure{$\widetilde G^\circ$ soit le sous-groupe de $\widetilde G$},} \inserted{Si $t\in\mathrm{Tr}(\widetilde X)$}, $\sigma_t$ échange $\tau$ et $\bar\tau$ \unsure{donc|dans} le groupe \inserted{$\widetilde G^\circ\subset\widetilde G$} \deleted{\unsure{engendré}} \inserted{formé par} par les nbrs pairs de symétries \deleted{\unsure{opère}} conserve chacune des $\tau\in H$ -- et les orbites de $\widetilde G^\circ$ sur $S-\widetilde X$ sont de cardinal $3$, et ce sont les $6$ triades de l'hexagone $H$.

\deleted{(Ainsi l'ens des trois \unsure{orb}} \deleted{d) Soit}

\textcircled{\scriptsize 11} \unsure{Soit|Soient} $(X_\alpha)_{\alpha\in A}$ une famille de trois bitriangles mutuellmt disjoints. Pour $\alpha\in A$, soit $I_\alpha$ l'ens de cardinal $3$ associé à $X_\alpha$ (\deleted{des} \inserted{l'ens} trois orbites de $\mathrm{AutInt}(X_\alpha)$ opérant sur $S-X_\alpha=\bigcup_{\beta\in A-\{\alpha\}}X_\beta$ On a donc une application $X_\beta\xrightarrow{p_{\alpha\beta}}I_\alpha$ $(\beta\ne\alpha)$
% archive sheet number: 50
